%%%%%%%%%%%%%%%%%%%%%%%%%%%%%%%%%%%%%%%%%%%%%%%%%%%%%%%%%%%%%%%%%%%%%%%%%%%%%%%%%%%%%%%%%%%%%%%%%%%%%%%%
\chapter{Models and Methods}
\label{ChapModels}

This chapter defines the common modelling and evaluation methodology used
throughout the gap-filling, forecasting and scenario-projection experiments.
Although these tasks use the same underlying NWFP data, they represent
different prediction settings. Gap-filling reconstructs observations removed
from within an historical record, forecasting predicts beyond a defined
forecast origin, and scenario projection evaluates a fitted model under
prescribed future environmental and management conditions. The corresponding
validation procedures are therefore described separately before introducing
the complete model roster, uncertainty-quantification framework and
interpretability method.

% This chapter collects every machine learning model, statistical baseline, uncertainty-quantification
% (UQ) method, and interpretability technique used anywhere in this dissertation, regardless of which
% downstream task (gap-filling, forecasting, or scenario projection) applies it. Chapters~\ref{Chap4}
% through~\ref{Chap6} assume the descriptions below and report only configuration, results, and
% task-specific discussion.

%%%%%%%%%%%%%%%%%%%%%%%%%%%%%%%%%%%%%%%%%%%%%%%%%%%%%%%%%%%%%%%%%%%%%%%%%%%%%%%%%%%%%%%%%%%%%%%%%%%%%%%%
\section{Evaluation Methods and Metrics}
\label{sec:evaluation_protocols}

The benchmark models optimise different internal objectives, including squared
error, absolute error, likelihood-based objectives and quantile loss.
Consequently, they were compared using common held-out evaluation metrics
rather than their model-specific training losses. Table~\ref{tab:task_protocols}
summarises the principal distinction between the three modelling tasks.

\begin{table}[htbp]
\centering
\small
\caption{Summary of the evaluation settings used in the dissertation.}
\label{tab:task_protocols}
\begin{tabular*}{\textwidth}{@{\extracolsep{\fill}}
>{\raggedright\arraybackslash}p{0.16\textwidth}
>{\raggedright\arraybackslash}p{0.28\textwidth}
>{\raggedright\arraybackslash}p{0.34\textwidth}
>{\raggedright\arraybackslash}p{0.14\textwidth}@{}}
\toprule
Task & Evaluation unit & Information available to the model & Primary metric \\
\midrule
Gap-filling &
Synthetic calendar gaps placed within the observed record &
Remaining observations and supporting variables before and after each gap &
$R^2$ \\
\addlinespace[0.5em]

Forecasting &
365-day prediction from five historical forecast origins &
Historical training data and realised supporting variables over the evaluation period &
Seasonal-mean-scaled MASE \\
\addlinespace[0.5em]

Scenario projection &
Conditional prediction under prescribed climate and management trajectories &
Historical training data and externally supplied scenario covariates &
Relative change from the reference scenario \\
\bottomrule
\end{tabular*}
\end{table}

\subsection{Gap-Filling Cross-Validation}

Gap-filling was evaluated as an interpolation problem over each tower's usable
measurement period. The evaluation domains were 1 October 2017--30 June 2019
for T2, 1 October 2017--31 December 2023 for T4, and
1 February 2020--31 December 2023 for T9. Only QC-valid observed
FCH\textsubscript{4} values within these domains were eligible for masking.

Five synthetic-gap scenarios were used: very short gaps of one hour, short
gaps of four hours, medium gaps of 32 hours, long gaps of 288 hours, and a
mixed scenario in which gap length was sampled from these four durations.
Non-overlapping calendar gaps were positioned randomly until approximately
25\% of the available target observations had been masked. Two independently
seeded repetitions were evaluated for each tower and gap scenario.

For every repetition, the masked FCH\textsubscript{4} observations were removed
before fitting the model. Supporting environmental and management variables
remained available, and the model could use historical context from either
side of a gap. This reflects the intended use of gap-filling, in which later
observations may be available when an earlier missing interval is
reconstructed. It is therefore distinct from prospective forecasting and
should not be interpreted as evidence of future predictive skill.

Metrics were first calculated independently for each repetition. The median
over repetitions was then calculated within each gap-duration scenario,
followed by the median across the five scenarios. This prevents a particularly
large or easy synthetic gap from dominating the headline result. Where a
combined value is reported across T2, T4 and T9, the tower-level results are
averaged with equal weight so that the longer T4 record does not automatically dominate
the comparison.


\subsection{Forecasting Cross-Validation}

Forecasting was evaluated through fixed-origin backtesting. Five forecast
origins were defined on 16 December in 2018, 2019, 2020, 2021 and 2022.
At each origin, a model was fitted using only data available on or before that
date and generated a continuous 365-day daily prediction. Tower--origin
combinations without sufficient historical or observed evaluation data were
excluded rather than assigned artificial scores.

The resulting forecasts were evaluated against QC-valid observed
FCH\textsubscript{4}, not against the reconstructed target series. The
gap-filled historical series could be used to initialise autoregressive
features, but predictions were scored only where a real observation was
available. This prevents agreement with the gap-filling model from being
mistaken for agreement with measured methane flux.

For recursive models, each predicted FCH\textsubscript{4} value was appended
to the available history and used to construct the autoregressive inputs for
subsequent days. No later observed methane value was inserted during the
365-day rollout. Direct tabular foundation-model configurations instead
predicted from the available covariate representation without requiring the
same recursive target-update mechanism.

Historical backtests supplied the realised supporting environmental variables
over each future evaluation window. The resulting experiment is therefore a
conditional methane-flux forecast: it evaluates the mapping from environmental
and management conditions to FCH\textsubscript{4}, together with any recursive
target dependence, rather than simultaneously evaluating a separate weather
forecast. Future climate variables are supplied by the climate-scenario data
in the scenario-projection experiment.

Performance was additionally divided into lead-time intervals of 1--7,
8--30, 31--90, 91--180, 181--270 and 271--365 days. This permits degradation
with forecast horizon to be distinguished from overall performance. Aggregate
metrics are weighted by the number of observed target values contributing to
each tower--origin--horizon cell, ensuring that cells containing very few real
observations do not receive the same influence as well-observed cells.

\subsection{Scenario-Projection Evaluation}

Scenario projection has no observed future target and is therefore treated as
a conditional projection rather than a forecast-validation exercise. The
best-performing forecasting configuration is held fixed across all prescribed
climate and management trajectories, which are compared with a common
reference over the same projection period.

For scenario \(s\), the relative change in mean predicted
FCH\textsubscript{4} is calculated as

\begin{equation}
\Delta_s
=
100
\left(
\frac{
\overline{\hat{y}}_s-\overline{\hat{y}}_{\mathrm{ref}}
}{
\overline{\hat{y}}_{\mathrm{ref}}
}
\right),
\end{equation}

where \(\overline{\hat{y}}_s\) and
\(\overline{\hat{y}}_{\mathrm{ref}}\) denote the mean predictions under the
scenario and reference trajectories, respectively. These values quantify
modelled conditional differences and are not interpreted as validated
forecasts of future methane flux.

\subsection{Prediction Metrics}

Let \(y_i\) denote an observed flux and \(\hat{y}_i\) the corresponding model
prediction for \(n\) valid evaluation observations. Mean absolute error and
root mean squared error were calculated as

\begin{equation}
\mathrm{MAE}
=
\frac{1}{n}
\sum_{i=1}^{n}
\left|y_i-\hat{y}_i\right|,
\end{equation}

and

\begin{equation}
\mathrm{RMSE}
=
\sqrt{
\frac{1}{n}
\sum_{i=1}^{n}
\left(y_i-\hat{y}_i\right)^2
}.
\end{equation}

MAE represents the typical absolute error in the original flux units, whereas
RMSE gives greater weight to the comparatively rare large errors associated
with methane peaks.

Gap-filling primarily follows the OLS \(R^2\) convention used by the published
NWFP comparator \cite{zhu2023a}. In the implemented evaluation, this is the
squared Pearson correlation between observations and predictions:

\begin{equation}
R^2 (OLS)
=
\left[
\frac{
\sum_{i=1}^{n}
(y_i-\bar{y})
(\hat{y}_i-\bar{\hat{y}})
}{
\sqrt{
\sum_{i=1}^{n}(y_i-\bar{y})^2
\sum_{i=1}^{n}(\hat{y}_i-\bar{\hat{y}})^2
}
}
\right]^2.
\label{eq:r2_ols}
\end{equation}

Unlike the residual coefficient of determination
\(1-\sum(y_i-\hat{y}_i)^2/\sum(y_i-\bar{y})^2\),
Equation~\ref{eq:r2_ols} is bounded between zero and one. It measures the
strength of linear agreement but is comparatively insensitive to additive
bias and scale error. MAE and RMSE are therefore retained as supporting
metrics so that a model cannot be assessed from correlation alone.

Forecasting uses a seasonal-mean-scaled form of mean absolute scaled error.
For each forecast origin, the reference prediction for a target date is the
mean observed FCH\textsubscript{4} for the corresponding day of year within a
circular \(\pm 7\)-day window, calculated exclusively from observations before
the forecast origin. The reported metric is

\begin{equation}
\mathrm{MASE}_{\mathrm{seasonal}}
=
\frac{
\mathrm{MAE}(y,\hat{y})
}{
\mathrm{MAE}(y,\hat{y}^{\,\mathrm{seasonal}})
}.
\label{eq:mase_seasonal}
\end{equation}

This applies the scale-free comparison principle of MASE
\cite{hyndman2006accuracy} using the historical seasonal mean as the explicit
reference forecast. Values below one indicate that the model improves on the
seasonal mean, a value of one indicates equal error, and values above one
indicate inferior performance.

Weighted absolute percentage error was calculated as

\begin{equation}
\mathrm{WAPE}
=
\frac{
\sum_{i=1}^{n}|y_i-\hat{y}_i|
}{
\sum_{i=1}^{n}|y_i|
}.
\end{equation}

WAPE is more stable than observation-wise percentage errors for this
application because FCH\textsubscript{4} is signed and frequently lies near
zero. Pearson correlation was also reported to describe agreement in temporal
variation independently of absolute calibration. No single metric is treated
as sufficient: MASE provides the primary forecasting comparison, while MAE,
RMSE, WAPE and correlation identify differences in scale, peak sensitivity and
temporal tracking.

\section{Models Used}
\label{sec:model_roster}

The benchmark intentionally includes simple baselines, conventional statistical
models, tree ensembles, locally trained neural architectures and pretrained
tabular foundation models. This breadth permits the claim that one model
family outperforms another to be tested empirically under common evaluation
conditions. Table~\ref{tab:model_roster} summarises how each model operates and
where it is used. ``GF'' and ``FC'' denote gap-filling and forecasting 
respectively.

% {
% \small
% \setlength{\tabcolsep}{4pt}
% \renewcommand{\arraystretch}{1.15}
% \begin{longtable}{
% @{}
% >{\raggedright\arraybackslash}p{0.17\textwidth}
% >{\raggedright\arraybackslash}p{0.60\textwidth}
% >{\raggedright\arraybackslash}p{0.15\textwidth}
% @{}
% }
% \caption{Model families and their roles in the dissertation.}
% \label{tab:model_roster}\\
% \toprule
% Model & Operating principle & Application \\
% \midrule
% \endfirsthead

% \multicolumn{3}{l}{\tablename\ \thetable\ -- continued}\\
% \toprule
% Model & Operating principle & Application \\
% \midrule
% \endhead

% \midrule
% \multicolumn{3}{r}{Continued on next page}\\
% \endfoot

% \bottomrule
% \endlastfoot

% Training-set mean &
% Predicts the arithmetic mean of the available training target for every
% held-out observation. It contains no temporal, environmental or management
% information and therefore provides a zero-information reference against which
% all fitted gap-filling models should improve. &
% GF \\

% Marginal Distribution Sampling &
% Fills a target from observations recorded under similar meteorological
% conditions \cite{wutzler2018reddyproc}. The implemented hierarchy first
% searches for observations with comparable shortwave radiation, air temperature
% and VPD; then relaxes the search to shortwave radiation; and finally uses an
% expanding mean-diurnal-course fallback. Unlike a fitted regression model, MDS
% constructs each estimate by averaging conditionally similar historical
% observations. &
% GF \\

% MICE &
% Uses multivariate imputation by chained equations. Each incomplete variable is
% treated in turn as a regression target conditional on the remaining variables,
% and the procedure cycles through the columns iteratively. The implementation
% uses Bayesian ridge regression as the fixed conditional estimator. Masked
% FCH\textsubscript{4} values are recovered from the jointly imputed
% feature--target matrix. &
% GF \\

% HyperImpute &
% Retains the chained-equations structure of MICE but replaces its single fixed
% conditional estimator with an AutoML search. A separate regression or
% classification algorithm may therefore be selected for each imputed column.
% This allows nonlinear and variable-specific relationships to be represented,
% at substantially greater computational cost. &
% GF \\

% Random Forest &
% Fits many regression trees to bootstrap samples of the training data while
% considering a random subset of predictors at each split. Individual trees
% capture nonlinear thresholds and interactions, while averaging their
% predictions reduces variance. A meteorology-only configuration reproduces the
% published baseline, whereas the expanded configuration includes environmental,
% calendar, livestock and management features. Cross-tower versions use tower
% indicators to distinguish the three measurement contexts. &
% GF, FC \\

% XGBoost &
% Constructs regression trees sequentially. Each new tree is fitted to improve
% the residual errors of the existing ensemble, while shrinkage, tree-depth
% constraints and explicit regularisation limit overfitting. Its sequential
% error-correction mechanism can model nonlinear feature interactions but may
% remain sensitive to sparse methane peaks. &
% GF, FC \\

% LightGBM &
% Also uses gradient-boosted decision trees, but relies on histogram-based split
% search and leaf-wise tree growth. Histogram binning reduces computational cost,
% while leaf-wise growth expands the branch offering the greatest loss
% reduction. Depth and leaf-size restrictions are used to control the resulting
% model complexity. &
% GF, FC \\

% Artificial Neural Network &
% Passes the input features through successive fully connected nonlinear layers
% and learns their weights through gradient-based optimisation. The ANN appears
% in the published-baseline replication rather than the final improved
% gap-filling roster, permitting comparison with earlier methane-flux studies. &
% GF replication \\

% Support Vector Machine &
% Uses support-vector regression to estimate a function within an
% \(\epsilon\)-insensitive error tube. Kernel transformations permit nonlinear
% relationships to be represented through similarities between observations,
% while the regularisation parameter controls the balance between a flatter
% function and errors outside the tube. It is retained only for published
% baseline replication. &
% GF replication \\

% SAITS &
% A self-attention architecture designed specifically for multivariate
% time-series imputation \cite{du2023saits}. Two diagonally masked
% self-attention blocks learn relationships between variables and time points
% without directly copying the value being reconstructed. Their estimates are
% combined through a learned weighting mechanism and trained using observed-value
% reconstruction and artificial-masking objectives. &
% GF \\

% LSTM and Bi-LSTM &
% Long short-term memory networks maintain a recurrent cell state controlled by
% input, forget and output gates. These gates regulate which historical
% information is stored, removed and exposed to the prediction layer. The
% bidirectional gap-filling model processes a sequence in both temporal
% directions, which is appropriate when observations after a historical gap are
% available. The forecasting LSTM is causal at inference and updates its input
% history using earlier predictions during recursive rollout. &
% GF, FC \\

% Temporal Fusion Transformer &
% Combines recurrent local processing with variable-selection networks, gated
% residual components and attention over the encoded sequence. Variable-selection
% weights allow the model to emphasise different observed and known-future
% inputs, while attention represents longer temporal relationships. The
% forecasting implementation uses a fixed historical lookback and produces a
% multi-step sequence, of which the next prediction is introduced into the
% recursive chain. &
% FC \\

% DLinear &
% Decomposes each input sequence into a moving-average trend and a residual
% seasonal component, then applies separate linear projections to the two
% components \cite{zeng2023}. The final forecast is their sum. It provides a
% deliberately simple neural forecasting baseline for testing whether elaborate
% attention or recurrent architectures are necessary. &
% FC \\

% SARIMAX &
% Represents the target through autoregressive terms, differencing and
% moving-average error terms, together with exogenous environmental and
% management regressors. Candidate autoregressive and moving-average orders are
% fitted separately at each tower and forecast origin, with the final order
% selected using AIC. SARIMAX provides a classical statistical comparison with
% the machine-learning models. &
% FC \\

% TabPFN &
% A pretrained tabular foundation model that learns a prediction procedure from
% large numbers of synthetic tabular tasks
% \cite{hollmann2023tabpfn,hollmann2025tabpfn}. At inference, labelled historical
% rows form an in-context training set and unlabelled rows form prediction
% queries. Transformer attention relates the queries to the supplied context,
% producing point and quantile predictions without conventional task-specific
% gradient training. For forecasting, temporal structure is supplied through
% calendar, lagged, environmental and management features, following the
% tabularisation principle used by TabPFN-TS \cite{hoo2026tabpfnts}. &
% GF, FC \\

% TabICLv2 &
% Extends tabular in-context learning to larger tabular contexts through revised
% synthetic pretraining and scalable attention
% \cite{qu2025tabicl,qu2026tabiclv2}. As with TabPFN, labelled context rows and
% query rows are processed jointly at inference. In gap-filling it is fitted
% separately by tower because its fixed context budget can be diluted by pooling.
% The scenario model also uses a per-tower direct-regression configuration with
% a restricted climate-compatible feature set. &
% GF, FC, SP \\

% Prediction ensemble &
% Combines predictions from independently fitted models. The forecasting
% benchmark evaluates an unweighted mean of Random Forest, XGBoost, LightGBM and
% SARIMAX, while the final TabPFN combination averages complementary TabPFN
% configurations. Averaging can reduce model-specific variance, although it can
% also dilute a stronger constituent with weaker members. &
% FC \\
% \end{longtable}
% }

  {
  \small
  \setlength{\tabcolsep}{4pt}
  \renewcommand{\arraystretch}{1.13}
  \begin{longtable}{
  @{}
  >{\raggedright\arraybackslash}p{0.17\textwidth}
  >{\raggedright\arraybackslash}p{0.60\textwidth}
  >{\raggedright\arraybackslash}p{0.15\textwidth}
  @{}
  }
  \caption{Model families and their roles in the dissertation.}
  \label{tab:model_roster}\\
  \toprule
  Model & Operating principle & Application \\
  \midrule
  \endfirsthead

  \multicolumn{3}{l}{\tablename\ \thetable\ -- continued}\\
  \toprule
  Model & Operating principle & Application \\
  \midrule
  \endhead

  \midrule
  \multicolumn{3}{r}{Continued on next page}\\
  \endfoot

  \bottomrule
  \endlastfoot

  Training-set mean &
  Predicts the mean observed training target for every held-out row, providing a
  zero-information reference. &
  GF \\

  MDS &
  Reconstructs flux from historical observations under similar meteorological
  conditions, followed by an expanding mean-diurnal-course fallback
  \cite{reichstein2005mds,wutzler2018reddyproc,zhu2023a}. &
  GF \\

  MICE &
  Iteratively predicts each incomplete variable from the remaining variables.
  The implementation uses Bayesian ridge regression as the conditional estimator
  \cite{vanbuuren2011mice}. &
  GF \\

  HyperImpute &
  Extends chained-equation imputation by using an AutoML search to select a
  conditional estimator separately for each variable
  \cite{jarrett2022hyperimpute}. &
  GF \\

  Random Forest &
  Averages regression trees fitted to bootstrap samples with random feature
  selection at each split, capturing nonlinearities while reducing individual-tree
  variance \cite{breiman2001randomforests,irvin2021,kim2020,zhu2023a}. &
  GF, FC \\

  XGBoost &
  Builds regularised regression trees sequentially, with each tree correcting
  residual errors left by the existing ensemble
  \cite{chen2016xgboost,irvin2021}. &
  GF, FC \\

  LightGBM &
  Uses histogram-based gradient boosting and leaf-wise tree growth to model
  nonlinear relationships efficiently \cite{ke2017lightgbm}. &
  GF, FC \\

  ANN &
  Learns nonlinear feature combinations through successive fully connected
  layers trained by gradient-based optimisation
  \cite{rumelhart1986backprop,kim2020}. &
  GF replication \\

  SVM &
  Uses support-vector regression to fit a regularised function within an
  $\epsilon$-insensitive error margin, with kernels permitting nonlinear
  relationships \cite{drucker1997svr,kim2020}. &
  GF replication \\

  SAITS &
  Uses diagonally masked self-attention blocks and learned combination weights
  to reconstruct multivariate time-series gaps \cite{du2023saits}. &
  GF \\

  LSTM and Bi-LSTM &
  Use gated recurrent memory to retain or discard temporal information.
  Bi-LSTM uses context on both sides of a historical gap, whereas forecasting
  LSTM operates causally during rollout
  \cite{hochreiter1997lstm,schuster1997bidirectional}. &
  GF, FC \\

  TFT &
  Combines recurrent processing, variable-selection networks, gating and
  temporal attention to represent local and longer-range forecasting
  relationships \cite{lim2021tft}. &
  FC \\

  DLinear &
  Decomposes a sequence into trend and residual components and forecasts each
  with a linear projection \cite{zeng2023}. &
  FC \\

  SARIMAX &
  Combines autoregressive, differencing and moving-average terms with exogenous
  environmental and management regressors. Candidate orders are selected using
  AIC \cite{box2015timeseries}. &
  FC \\

  TabPFN &
  Uses a transformer pretrained on synthetic tabular tasks to infer predictions
  from labelled context rows and unlabelled query rows
  \cite{hollmann2023tabpfn,hollmann2025tabpfn}. Temporal prediction is represented
  through engineered time-series features \cite{hoo2026tabpfnts}. &
  GF, FC \\

  TabICLv2 &
  Uses scalable tabular in-context learning with synthetic pretraining
  \cite{qu2025tabicl,qu2026tabiclv2}. &
  GF, FC \\

  Prediction ensemble &
  Averages predictions from complementary fitted models to reduce
  model-specific variance. Both tree--statistical and TabPFN combinations are
  evaluated \cite{bates1969combination}. &
  FC \\

  \end{longtable}
  }

For locally fitted tabular models, missing predictor values were replaced using
statistics estimated from the corresponding training data. Pooled models used
tower indicators to share information while retaining tower-specific effects.
Exact feature configurations, context lengths and hyperparameters are reported
with the corresponding experiments; Table~\ref{tab:model_roster} defines only 
the common operating principles used throughout the results chapters.

% For locally fitted tabular models, missing predictor values were replaced using
% statistics estimated from the corresponding training data. Where models were
% pooled across towers, tower-indicator variables allowed a common fitted
% function to retain tower identity. This differs from treating the three towers
% as interchangeable observations: pooling shares statistical information while
% still permitting tower-specific offsets and interactions.

% Exact feature configurations, context lengths and hyperparameters are reported
% with the corresponding gap-filling, forecasting and scenario-projection
% experiments. The purpose of Table~\ref{tab:model_roster} is to define the
% operating principles of the models once, avoiding repeated generic
% descriptions in each results chapter.



%%%%%%%%%%%%%%%%%%%%%%%%%%%%%%%%%%%%%%%%%%%%%%%%%%%%%%%%%%%%%%%%%%%%%%%%%%%%%%%%%%%%%%%%%%%%%%%%%%%%%%%%
\section{Uncertainty Quantification}
\label{sec:uncertainty_methods}

Point predictions alone do not distinguish confident predictions from periods
in which multiple methane outcomes are plausible. Uncertainty was therefore
represented through quantile-based machine learning and conformal calibration.
For a conditional target distribution \(Y\mid X=x\), the quantile model
estimates

\begin{equation}
Q_{\tau}(x)
=
\inf\left\{
y:
P(Y\leq y\mid X=x)\geq\tau
\right\},
\end{equation}

where \(\tau\in(0,1)\) denotes the desired quantile. The experiments use
\(\tau=0.05\), \(0.50\) and \(0.95\), producing a median point prediction and a
nominal 90\% interval.

Quantile models are trained or evaluated using pinball loss
\cite{koenker1978quantiles}:

\begin{equation}
L_{\tau}(y,\hat{q}_{\tau})
=
\begin{cases}
\tau(y-\hat{q}_{\tau}),
& y\geq\hat{q}_{\tau},\\
(1-\tau)(\hat{q}_{\tau}-y),
& y<\hat{q}_{\tau}.
\end{cases}
\label{eq:pinball}
\end{equation}

The asymmetric weighting means that underprediction is penalised more heavily
for a high quantile and overprediction more heavily for a low quantile.
Minimising Equation~\ref{eq:pinball} therefore estimates the corresponding
conditional quantile rather than the conditional mean.

% Several mechanisms were used to obtain initial intervals. Quantile Regression
% Forest derives quantiles from the distribution of training outcomes represented
% within the fitted forest leaves. XGBoost and LightGBM fit separate
% quantile-objective models for the 0.05, 0.50 and 0.95 levels. SARIMAX obtains
% intervals from its fitted state-space forecast distribution, while the TFT
% uses a quantile-output head. TabPFN and TabICLv2 provide native predictive
% quantiles from their pretrained regression distributions. Where an ensemble
% was evaluated, corresponding lower, median and upper predictions from its
% constituent models were combined using the same weights as the point ensemble.

Several mechanisms were used to obtain initial prediction intervals. Quantile
Regression Forest estimates conditional quantiles from the distribution of
training outcomes represented within the fitted forest leaves
\cite{meinshausen2006qrf}. XGBoost and LightGBM fit separate
quantile-objective models for the 0.05, 0.50 and 0.95 levels, while SARIMAX
obtains intervals from its fitted state-space forecast distribution. The TFT
uses a quantile-output head, and TabPFN and TabICLv2 provide predictive
quantiles from their pretrained regression distributions. For an ensemble,
corresponding lower, median and upper predictions were combined using the same
weights as the point predictions.

Raw model quantiles are not guaranteed to achieve their nominal coverage on
NWFP data. Split-conformal calibration was therefore used to adjust the
intervals using held-out residual information
\cite{lei2018distributionfree}.
For symmetric split-conformal intervals, the calibration score for observation
\(i\) is

\begin{equation}
s_i
=
\left|
y_i-\hat{q}_{0.50}(x_i)
\right|.
\end{equation}

For a desired miscoverage rate \(\alpha=0.10\), the conformal margin is the
finite-sample corrected \(1-\alpha\) quantile of the calibration scores. The
calibrated interval is then

\begin{equation}
C_{0.90}(x)
=
\left[
\hat{q}_{0.50}(x)-\hat{s},
\hat{q}_{0.50}(x)+\hat{s}
\right].
\end{equation}

Forecasting calibration used a leave-one-anchor-out procedure. When one
forecast origin was evaluated, residuals from the remaining four origins were
pooled within the corresponding lead-time interval. The resulting conformal
margin was then applied to the held-out origin. This preserves separation
between calibration and evaluation while allowing interval width to vary
between short and long forecast horizons.

Conformalized Quantile Regression was additionally used where the native
quantile spread contained information that a symmetric interval around the
median would discard \cite{romano2019cqr}. Its nonconformity score is

\begin{equation}
s_i^{\mathrm{CQR}}
=
\max\left\{
\hat{q}_{0.05}(x_i)-y_i,\;
y_i-\hat{q}_{0.95}(x_i)
\right\},
\end{equation}

giving the calibrated interval

\begin{equation}
C^{\mathrm{CQR}}_{0.90}(x)
=
\left[
\hat{q}_{0.05}(x)-\hat{s}_{\mathrm{CQR}},
\hat{q}_{0.95}(x)+\hat{s}_{\mathrm{CQR}}
\right].
\end{equation}

For gap-filling, QRF and TabICLv2 intervals were evaluated on the synthetic
held-out gaps. Calibration data were separated from the interval-evaluation
data within each tower and gap-duration setting. For forecasting, calibration
was separated by forecast origin and lead-time interval. Scenario intervals
reuse margins learned from historical backtests; they quantify historically
calibrated predictive dispersion but cannot guarantee coverage under an
unobserved future distribution shift.

Let \(\hat{l}_i\) and \(\hat{u}_i\) denote the final lower and upper
bounds of the interval evaluated for observation \(i\). These may correspond
to either the raw model quantiles or their conformally calibrated
counterparts.

Three complementary interval metrics were reported. Prediction Interval
Coverage Probability is

\begin{equation}
\mathrm{PICP}
=
\frac{1}{n}
\sum_{i=1}^{n}
\mathbb{I}
\left(
\hat{l}_i
\leq y_i
\leq
\hat{u}_i
\right).
\end{equation}

This measures the proportion of observations inside the interval. For a nominal
90\% interval, values substantially below 0.90 indicate undercoverage and
overconfidence, whereas values well above 0.90 indicate conservative
intervals.

Mean Prediction Interval Width is

\begin{equation}
\mathrm{MPIW}
=
\frac{1}{n}
\sum_{i=1}^{n}
\left(
\hat{u}_i-\hat{l}_i
\right).
\end{equation}

MPIW measures sharpness. Narrower intervals are preferable only when their
coverage remains comparable; an arbitrarily narrow interval with poor PICP is
not useful. Mean pinball loss across the three quantiles measures the accuracy
of the quantile predictions. PICP, MPIW and pinball loss are
therefore interpreted jointly as coverage, sharpness and quantile accuracy,
respectively.

% ============================================================================
\section{Interpretability}
\label{sec:interpretability_methods}

Permutation-based prediction sensitivity was used as the common model-agnostic
interpretability method. For each forecast origin and tower, the fitted model
first produced an unmodified prediction vector
\(\hat{\mathbf{y}}\). Predictor \(j\) was then shuffled across the evaluation
rows and the already fitted model was applied again without retraining.

Sensitivity was calculated as

\begin{equation}
I_j
=
\left|
\frac{1}{n}
\sum_{i=1}^{n}
\hat{y}^{\,\pi(j)}_i
-
\frac{1}{n}
\sum_{i=1}^{n}
\hat{y}_i
\right|,
\label{eq:permutation_importance}
\end{equation}

where \(\hat{y}^{\,\pi(j)}_i\) is the prediction obtained after permuting
feature \(j\). The score therefore measures the change in mean predicted
FCH\textsubscript{4} caused by disrupting the information carried by that
feature. Unlike conventional error-based permutation importance
\cite{breiman2001randomforests}, it does not require an observed future target
and can consequently be applied to conditional scenario projections.

Permutations were performed within each tower so that values were not mixed
between T2, T4 and T9. Tower indicators were not permuted. Scores were
calculated across the five forecast origins using deterministic seeds and were
then summarised overall and by tower.

These scores measure model reliance rather than physical causality. Correlated
predictors may share importance, while permutation may create combinations of
variables that are uncommon in the observed data. Interpretation therefore
focuses on stable feature groups and cross-tower patterns rather than small
differences between individual rankings.
